\clearpage
\fancyhead[L]{\small IRIS.}
\SourceLeaf{479}
\begin{center}
\rule{0.55\linewidth}{0.4pt}\\[2pt]
Astronomische Nachrichten. Band XXVI. Nr. 611. Seite 173. 174. 1847 Sept. 16.\\[2pt]
\rule{0.55\linewidth}{0.4pt}
\end{center}
\begin{center}[Beobachtungen am Meridiankreise von \textsc{Gauss}.]\end{center}

\begin{center}
\small
\setlength{\tabcolsep}{3pt}
\begin{tabular}{lrll}
Iris: 1847 Aug. 21. & $9^{\rm u}\ 52^{\rm m}\ 54^{\rm s}1$ Gött. m. Z. & $297^\circ\ 47'\ 20''8$ & $-13^\circ\ 42'\ 20''4$ \\
\end{tabular}
\end{center}

Die Antritte an den einzelnen Fäden stimmten nicht besonders, da wegen ziehender Wolken das Gestirn schwer zu sehen war. Unmittelbar darauf dehnten sich die Wolken noch weiter aus, so dass die Culmination von $\alpha$ Steinbock nicht beobachtet werden konnte.

\begin{center}
\small
\setlength{\tabcolsep}{4pt}
\begin{tabular}{lrrr}
Neptun 1847 Aug. 17. & $12^{\rm u}\ 22^{\rm m}\ 55^{\rm s}2$ & $331^\circ\ 27'\ 14''1$ & $-12^\circ\ 21'\ 5''0$ \\
\end{tabular}
\end{center}

\clearpage
\hypertarget{native-479-lower}{}
\begin{center}
\rule{0.55\linewidth}{0.4pt}\\[2pt]
Astronomische Nachrichten. Band XXVI. Nr. 613. Seite 193..208. 1847 October 14.\\[2pt]
\rule{0.55\linewidth}{0.4pt}
\end{center}
\begin{center}[Beobachtungen der \emph{Iris} am Meridiankreis von \textsc{Gauss}.]\end{center}
\begin{flushright}Göttingen 1847 Sept. 5.\end{flushright}

Herr \textsc{Gould} hat mir eben die aus meinen Beobb. der Iris vom 21. 28. August, in Verbindung mit der Cambridger vom 14. abgeleiteten Elemente geschickt, die ich Ihnen mitzutheilen eile. --- --- --- ---

Gestern Abend habe ich Iris noch am Meridiankreise beobachtet.

\begin{center}
\setlength{\tabcolsep}{4pt}
\begin{tabular}{l|r|r|r}
 & Göttinger M. Z. & AR. & Declination \\\hline
1847 Sept. 4 & $8^{\rm u}\ 51^{\rm m}\ 12^{\rm s}0$ & $296^\circ\ 7'\ 14''1$ & $-14^\circ\ 6'\ 3''8$ \\
\hfill 6 & $8\ 42\ 54.5$ & $296\ 0\ 46.5$ & $-14\ 8\ 56.1$ \\
\hfill 10 & $8\ 26\ 43.6$ & $295\ 53\ 59.8$ & $-14\ 13\ 59.5$ \\
\end{tabular}
\end{center}

Der Planet erschien bei dieser Beobachtung lichtschwächer als bei irgend einer früheren. Die Fäden blieben ganz undeutlich, wenn der Planet sichtbar sein sollte. Viel mag an dem Zustande der Luft gelegen haben, denn alle Sterne erschienen (dem blossen Auge), während der Himmel ganz wolkenfrei war, in mattem nebelichten Lichte.

\begin{center}
\begin{tabular}{l|r|r|r}
Sept. 12 & $8^{\rm u}\ 18^{\rm m}\ 50^{\rm s}5$ & $295^\circ\ 53'\ 37''8$ & $-14^\circ\ 16'\ 16''3$ \\
\end{tabular}
\end{center}

Din Luft war an diesem Abend günstiger als am 10. Dagegen war am 13. die Iris nur im ganz verdunkelten Felde eben zu sehen, und eine Beobachtung unmöglich. Der um etwa $1^{\rm m}\ 40^{\rm s}$ vorangehende auf fast gleichem Parallel stehende Stern 7. 8. Grösse hatte sich noch bei nothdürftiger Beleuchtung beobachten lassen.

\begin{center}
\begin{tabular}{l|r|r|r}
Sept. 19 & $7^{\rm u}\ 52^{\rm m}\ 19^{\rm s}9$ & $296^\circ\ 8'\ 51''1$ & $-14^\circ\ 22'\ 30''0$ \\
\end{tabular}
\end{center}
